% Source: https://artofproblemsolving.com/community/c3h1761574_new_years_problems_ii
% Stable identity: aops-topic-1761574-q03
First count arrangements on ten labelled chairs, then divide by $10$: since all people are distinct, no rotation other than the identity fixes an arrangement, so each class of rotations has exactly $10$ members.

\textbf{Twins:} $10$ pairs of adjacent chairs and $2$ orders: $20$ ways. By symmetry, say they occupy chairs $0$ and $1$.

\textbf{Jessica and John} use two of the chairs $2,3,\dots,9$, which form a row (chairs $9$ and $2$ are separated by the twins' chairs). Pairs at distance $1$: $(2,3),\dots,(8,9)$, $7$ pairs; at distance $2$: $(2,4),\dots,(7,9)$, $6$ pairs (chairs $9$ and $1$, or $0$ and $2$, would involve a twin). So $\dbinom82-13=15$ unordered pairs are allowed, and $30$ ordered placements.

\textbf{Adults:} choose $3$ of the remaining $6$ chairs, $\dbinom63=20$ ways; of the $6$ ways to seat the adults in them, exactly $3$ give the clockwise order Mr.\ Smith, Mrs.\ Smith, Annie (one cyclic order out of two). So $60$ ways.

With labelled chairs: $20\cdot30\cdot60=36000$. Up to rotation: $\dfrac{36000}{10}=\boxed{3600}$.
